\chapter{Conclusion and Future Work}
\label{ch:conclusion}

This final chapter summarizes the contributions of the project
against the five objectives defined in Chapter~\ref{ch:introduction},
acknowledges the limitations of the current system, and outlines
four directions for future work.

\section{Summary of Contributions}
\label{sec:conclusion-contrib}

The project as presented in this report has delivered five concrete
contributions, each aligned with one of the five objectives from
Section~\ref{sec:objectives}:

\begin{enumerate}
  \item \textbf{A four-class driver stress classifier based on a
  bidirectional LSTM with additive attention.} The classifier
  operates on a five-feature, 30-timestep window sampled at 4~Hz
  and uses only wrist-accessible signals. Its design follows the
  bidirectional-LSTM blueprint of \citet{graves2005bilstm} and the
  attention intuition of \citet{vaswani2017attention}, scaled down
  to a 160{,}000-parameter compact form factor.

  \item \textbf{Empirical accuracy of 97.8\% and weighted F1 of 0.978}
  on the three dominant classes of the held-out WESAD test split
  \citep{schmidt2018wesad}, with a deployed TensorFlow Lite
  \citep{tflite2023} footprint of 72~KB. This exceeds both the accuracy
  target ($\ge 90\%$) and the size target ($\le 80$~KB) set in
  Chapter~\ref{ch:introduction}.

  \item \textbf{A Wear~OS Kotlin application that performs real-time
  on-device inference.} The application collects wrist sensors at
  their native rates, resamples them to a common 4~Hz grid, applies
  the cached per-feature normalization learned at training time,
  runs the quantized TensorFlow Lite interpreter, and emits a
  smoothed stress estimate every thirty seconds with end-to-end
  latency well below the 50~ms-per-window target.

  \item \textbf{A documented twelve-byte BLE GATT protocol}
  (service \texttt{ABCD1234\ldots}, characteristic
  \texttt{ABCD1235\ldots}) that carries a timestamp, a 4-valued
  stress level, and a confidence. The protocol is fully specified
  in Section~\ref{sec:m-ble} and can be implemented by any BLE
  central or peripheral independently of the reference
  implementation.

  \item \textbf{A graduated, physiology-driven adaptive cabin
  control system} implemented on the Arduino Nano 33 BLE, comprising
  four active cabin profiles (Calm, Mild, Moderate, Critical), a
  safe \texttt{WAITING} state, and a sixty-second BLE timeout.
  The fail-safe behavior mirrors the design philosophy of the
  AlcoWatch BAC system under provisional patent ACN1408
  \citep{acn1408}.
\end{enumerate}

These five contributions add up to a working end-to-end prototype.
A driver wearing a Wear~OS device with this application installed,
while operating a vehicle fitted with the reference Arduino module,
will see their cabin lighting, ventilation, and audio adapt
automatically to their autonomic-nervous-system state, with no cloud
dependency and no compromise of sensor data privacy.

\section{Limitations}
\label{sec:conclusion-limits}

Four limitations must be acknowledged. Addressing them is what the
future-work program of Section~\ref{sec:conclusion-future} targets.

\textbf{Laboratory versus naturalistic stressors.} WESAD
\citep{schmidt2018wesad} was collected under controlled laboratory
conditions: subjects were stationary, posture was constrained, and
stressors were delivered via a fixed protocol (a Trier Social Stress
Test variant). Real driving stress is different. It mixes with
physical workload (steering, pedals), sustained mental demand, and
ambient vibration. The physiological signatures might differ enough
to reduce the classifier's accuracy in actual deployment, which is
an open empirical question.

The second limitation is \textbf{lack of real-world deployment testing}.
The system has not been evaluated on real drivers under real driving
conditions. The evaluation here is limited to held-out WESAD data and
scripted BLE simulations. A supervised on-road study is needed before
any safety claim can be made.

\textbf{Synthetic Critical class.} The ``Critical'' class was defined
using a per-subject quantile split of the WESAD stress condition, not
from a real driving emergency. That's an approximation of the autonomic
response a genuine near-miss would produce. Whether it's a close enough
approximation for the Critical cabin profile to fire at the right
moments is something only real-world testing can answer.

\textbf{Accelerometer magnitude discards direction.} The model uses
the three-axis accelerometer magnitude as a single scalar feature.
That keeps the feature count at five, but throws away directional
information. For example, a sharp lateral swerve and a vertical
pothole impact produce similar magnitudes but very different risk
profiles, and the model can't tell them apart. Admitting the three
raw axes as separate features would grow the input to seven and
slightly increase model size, but would likely improve discrimination.

\section{Future Work}
\label{sec:conclusion-future}

Four directions suggest themselves for follow-on work.

\textbf{Unification with AlcoWatch into a single platform.}
The AlcoWatch BAC system \citep{acn1408} and this stress-detection
system share most of their architecture: the same wrist sensors,
the same on-device TensorFlow Lite runtime, the same Arduino-based
vehicle controller, and the same fail-safe default-deny philosophy.
Merging them into a single Wear~OS application that runs both
inference engines concurrently and exposes two BLE characteristics
on one shared service is a natural next step. The Arduino firmware
would combine both decision logics: BAC-based ignition blocking
(AlcoWatch-style) and stress-driven cabin adaptation (this work).
That would spread the hardware cost of the wearable across two
distinct safety outcomes.

A second priority is \textbf{real-world data collection during
driving}. A follow-on project should collect a WESAD-style multimodal
dataset under naturalistic driving conditions, either on a simulator
or on a test track (subject to ethical approval). That dataset would
let the classifier be fine-tuned on the exact signal distribution it
encounters in deployment and would provide the empirical basis for a
quantitative safety claim.

\textbf{Personalization via per-driver calibration} is the third
direction. The per-subject quantile logic in the stress-split label
mapping (Section~\ref{sec:m-labels}) is evidence that between-subject
differences in baseline autonomic tone are large enough to matter.
A lightweight calibration step -- for example, a five-minute baseline
capture the first time the watch is worn, with the resulting per-user
statistics used to normalize incoming sensor streams -- would likely
improve Critical-class recall on drivers whose resting EDA or heart
rate sits at the extremes of the population.

Finally, \textbf{knowledge distillation to a sub-25~KB student model}
would reduce inference latency and energy consumption and would open
the door to ultra-low-power wearables and hearables. The current
72~KB model runs fine on Wear~OS, but a smaller model would be useful
for wristband-class devices. One practical route is knowledge
distillation: train a student network (a single-direction LSTM or a
1-D CNN with under 20{,}000 parameters) to match the soft predictions
of the current BiLSTM+attention teacher. If the student can hold 90\%
accuracy in a 25~KB footprint, continuous always-on stress monitoring
becomes feasible on much smaller hardware.

Each of these directions extends rather than replaces the contributions
of this report. Together they sketch the path from a validated
laboratory prototype to a deployable, personalized, multi-purpose
wearable safety platform.
